\chapter{隐藏状态估计：观测算子、可辨识性与滤波}
\label{ch:est}

全场 $\vect{X}$ 不能被家电传感器直接读取。本章规定：哪些量允许进入估计器、逆问题为什么病态、两节点参数在什么激励下可分离，以及在线用哪一种滤波器。估计器的输出是决策状态 $\vect{z}$ 与参数 $\vect{\theta}$ 的后验，不是一个温度点。

\section{可安装观测}

观测向量分成食品侧与安全侧。安全侧不参与品质滤波，只进入第 \ref{ch:safe} 章的硬约束。

\begin{table}[H]
\centering
\caption{进入品质估计器的观测}
\begin{tabularx}{\textwidth}{llX}
\toprule
通道 & 符号 & 与隐藏状态的关系 \\
\midrule
腔体空气温度 & $y_T$ & 靠近 $T_\infty$，受加热器辐射污染，需与壁温分开标定 \\
腔体相对湿度 & $y_\varphi$ & 由 $p_v/p_{\mathrm{sat}}(T_\infty)$ 决定，饱和时丢失分辨率 \\
食品质量 & $y_m$ & $m(t)=m_0-\int_\Gamma \dot m_\Gamma\,\mathrm{d}A$，直接约束累计蒸发 \\
表面热像 & $\vect{y}_{\mathrm{IR}}$ & $\epsilon(M_s,S)\,T_s^4$ 的空间采样，发射率本身未知 \\
电功率 & $y_P$ & 已知输入，写入边界，不当作食品状态 \\
排气温度 & $y_{T,e}$ & 冷凝与漏汽的积分证据 \\
\bottomrule
\end{tabularx}
\end{table}

\begin{definition}[观测方程]
\begin{equation}
\label{eq:obs}
\vect{y}_k=h\big(\Phi(\vect{X}(t_k)),\vect{\theta}\big)+\vect{v}_k,
\qquad \vect{v}_k\sim\mathcal{N}(\vect{0},\vect{R}_k).
\end{equation}
$h$ 至少包括：腔温一阶惯性、湿度传感器的饱和函数、质量的积分、红外亮度
\begin{equation}
\label{eq:ir}
L_j=\epsilon(\varepsilon_{l,s},S_s)\,\sigma T(\vect{r}_j)^4
+\big(1-\epsilon(\varepsilon_{l,s},S_s)\big)L_j^{\mathrm{refl}}.
\end{equation}
反射项 $L^{\mathrm{refl}}$ 来自加热管与腔壁。忽略它会把管壁辐射误认成食品表面温度。
\end{definition}

发射率 $\epsilon$ 随表面变干从湿表面的高值下降到结壳后的较低值。灵敏度分析扫描文献给出的 $\epsilon$ 区间，观测噪声用区间端点分别滤波。系统精度取该扫描下 $\vect{z}$ 分位区间的包络，不取某一热像器件的标称精度。

\section{逆问题的病态结构}

由表面温度反推内部场，是热传导逆问题。考虑厚度方向的线性化热方程，初始温度未知、边界热流未知时，解对高频测量噪声按 $\mathrm{e}^{kL}$ 放大，$k$ 为空间频率，$L$ 为特征厚度。因此：

\begin{enumerate}
\item 只加密表面红外像素，不能稳定恢复冷点；
\item 必须有一个穿透时间尺度上的动态激励，或一个对内部热阻敏感的积分量；
\item 质量损失对累计蒸发敏感，对冷点位置不敏感，不能单独代替温度场；
\item 先验（食材类别、厚度、反应动力学）把解限制在决策流形附近，正则化项来自物理，而不是一个随意的光滑惩罚。
\end{enumerate}

\begin{proposition}[静态观测不足以分离内外热阻]
在 \eqref{eq:twonode} 中令 $\dot T_1=\dot T_2=0$ 且 $\dot m=0$。稳态只给出
\begin{equation}
T_1=T_2=T_\infty
\end{equation}
（无蒸发时）。$R_{\mathrm{int}}$、$C_1$、$C_2$ 全部从稳态方程中消失。持续恒温加热的腔温读数不包含内部热阻信息。
\end{proposition}

这是“到时按菜单停机”失败的数学形式：菜单在稳态配方上标定，真实食品的 $R_{\mathrm{int}}/C_2$ 从未被这次烹饪观察过。

\section{动态可辨识性}

把 \eqref{eq:twonode} 在无相变工作点线性化，状态 $(\delta T_1,\delta T_2)$，输入 $(\delta T_\infty,\delta v)$。风速只进入 $R_{\mathrm{ext}}$。可观测矩阵对输出 $y=T_1$ 的秩为 2，内部状态在理论上可观测，但 Gram 矩阵的条件数随 $R_{\mathrm{int}}/R_{\mathrm{ext}}$ 增大而变坏。厚、干、结壳的食品，表面温度对冷点的响应又慢又弱。

相变项把输入通道从“只改 $T_\infty$”扩展为“改 $T_\infty$ 与 $q_{\mathrm{lat}}$”。短蒸汽脉冲使 $q_{\mathrm{lat}}$ 由负变正，表面节点出现一个与 $R_{\mathrm{ext}}$ 不同方向的跃迁。线性化输入矩阵变为两列，Fisher 信息矩阵对 $(R_{\mathrm{int}},C_2,k_m)$ 的最小特征值上升。局部灵敏度满足
\begin{equation}
\label{eq:sens}
\frac{\mathrm{d}}{\mathrm{d}t}\frac{\partial \vect{s}}{\partial\theta_i}
=\vect{A}(\vect{\theta})\frac{\partial \vect{s}}{\partial\theta_i}
+\frac{\partial \vect{A}}{\partial\theta_i}\vect{s}
+\frac{\partial \vect{B}}{\partial\theta_i}\vect{u},
\end{equation}
其中 $\vect{s}=(T_1,T_2,m)$。Fisher 信息
\begin{equation}
\label{eq:fisher}
\vect{\mathcal{I}}(\vect{\theta};\vect{u}_{0:K})
=\sum_{k=0}^{K}
\Big(\frac{\partial h}{\partial \vect{s}}\frac{\partial \vect{s}_k}{\partial\vect{\theta}}\Big)^{\!\top}
\vect{R}_k^{-1}
\Big(\frac{\partial h}{\partial \vect{s}}\frac{\partial \vect{s}_k}{\partial\vect{\theta}}\Big).
\end{equation}
探针设计在第 \ref{ch:ati} 章最大化 $\log\det\vect{\mathcal{I}}$ 的增量。

\section{贝叶斯估计器}

在线采用分层估计，避免把网格场和标量参数放进同一个高维粒子滤波器。

\begin{enumerate}
\item \textbf{参数层。} $\vect{\theta}=(C_1,C_2,R_{\mathrm{int}},R_{\mathrm{ext}},k_m,K_{\mathrm{ev}})$。先验按食材类别从知识库读取。更新用无迹卡尔曼或低维粒子滤波，过程噪声只加在会漂移的项（结壳使 $R_{\mathrm{int}}$ 随 $S$ 增大）。
\item \textbf{场层。} 一维径向降阶模型的模态系数用集合卡尔曼滤波（EnKF）更新。集合成员共享当前 $\vect{\theta}$ 后验的抽样，观测只含 $y_m$、$\vect{y}_{\mathrm{IR}}$ 与 $y_T$。
\item \textbf{决策层。} 每个集合成员计算 $\vect{z}^{(i)}=\Phi(\vect{X}^{(i)})$，经验分布即为 $p(\vect{z}\mid\vect{y}_{0:k})$。控制器使用分位数，不使用集合均值一个数。
\end{enumerate}

预测步是离散化的 \eqref{eq:twonode} 或径向格式。分析步是集合卡尔曼更新：预测协方差加膨胀因子 $\beta\ge 1$，空间局部化用 Gaspari--Cohn 函数，半径取一至两个径向节点。实时集合规模 $N=32$，离线对照 $N=128$。膨胀不足时用第 \ref{ch:ati} 章的探针残差发现，不在分析步里人为缩小观测噪声。

红外通道的似然使用 \eqref{eq:ir}，并把 $\epsilon$ 作为慢变参数与 $S_s$ 绑定。湿度通道在 $y_\varphi>0.98$ 时自动增大 $R_k$ 的对应对角元，表示饱和段几乎没有信息。

\section{可观测性失效时的行为}

出现下列任一情况，估计器不输出“点预测停机”，只输出不确定度，并把决策权交给主动探针或安全保持：

\begin{itemize}
\item $\vect{\mathcal{I}}$ 的最小特征值低于类别阈值（厚件、无质量信号、热像被蒸汽遮挡）；
\item 质量通道与湿度通道互相矛盾，残差超过 $3\sigma$ 持续一个热时间常数；
\item 集合成员的 $\alpha_m$ 分位区间同时覆盖“未达安全”和“已经过熟”。
\end{itemize}

第三种情况是蒸不熟与过蒸之间的决策冲突，必须由探针或延长低功率保温来压缩区间，不能靠把阈值调低来“选边”。

\section{本章与前后章的接口}

本章产出 $p(\vect{z},\vect{\theta}\mid\vect{y}_{0:k})$ 与 $\vect{\mathcal{I}}$。第 \ref{ch:ati} 章读取它们决定要不要探针；第 \ref{ch:causal} 章读取结构方程的当前参数；第 \ref{ch:safe} 章读取 $\alpha_m$ 与 $T$ 的分位数做机会约束。任何下游模块若只接收均值，视作违反数据合同。

观测标定、一维病态算例、Fisher 矩阵的数值秩和集合更新的完整伪代码见附录。退化后的输出只有区间，没有点估计停机时刻。
